\section{Related Work}

\textbf{Paper 10 baseline.} Paper~10~\cite{paper10adaptive} introduced
the one-threshold magnitude detector and evaluated it at the single
$\tau = 0.05$. This paper inherits the procedure and sweeps $\tau$.

\textbf{The Papers 7--10 sequence.}
Paper~7~\cite{paper7online} identified the K=200 hazard;
Paper~8~\cite{paper8multi} ruled out smoothing as a fix;
Paper~9~\cite{paper9selftraj} re-attributed Paper~6's K=200 collapse
to selection-policy coupling;
Paper~10 introduced the adaptive detector. This paper completes the
arc by characterising the $\tau$ axis the detector left implicit.

\textbf{Pareto-frontier analysis.} Pareto curves over a single tuning
parameter are standard in detection theory (ROC curves) and
forecasting (bias-variance trade-offs). The $\tau$-sweep here is the
analog in the calibration-strategy space.

\textbf{Change-point detection.} CUSUM~\cite{page1954cusum} and
Bayesian online change-point methods would each induce different
Pareto curves than the simple magnitude detector. Whether either
shifts the H3 operating point into the feasible region is future
work.

\textbf{Multi-paper sequence.} HygieneBench (Paper~3)~\cite{paper3hygienebench},
HygienePrio (Paper~4)~\cite{paper4hygieneprio}, Temporal Stability
(Paper~5)~\cite{paper5temporal}, Capacity-Indexed Decay
(Paper~6)~\cite{paper6capacity}, Online Calibration (Paper~7),
Multi-History Calibration (Paper~8), Self-Trajectory
(Paper~9), and Adaptive Recalibration (Paper~10) form the
methodological substrate.

\textbf{Policy mandates.} CISA BOD~22-01~\cite{cisa2021bod2201},
23-01~\cite{cisa2022bod2301}, NIST SP~800-40~\cite{nist2022sp80040r4},
2026 DBIR~\cite{verizon2026dbir}.
